\documentclass[a4paper]{report}
% Shared preamble for the English probability and statistics documents.
% Based on the original course preamble. Compile with pdfLaTeX.
\usepackage[top=1.5cm,bottom=2cm,left=2cm,right=2cm]{geometry}
\usepackage[T1]{fontenc}
\usepackage{lmodern}
\usepackage[utf8]{inputenc}
\usepackage[english]{babel}

\usepackage{amsmath}
\usepackage{amssymb}
\usepackage{amsthm}
\usepackage{mathrsfs}
\usepackage{mathtools}
\usepackage{etoolbox}
\usepackage{tikz}
\usetikzlibrary{shapes.multipart,arrows.meta}
\usepackage{xcolor}
\usepackage{enumerate}
\usepackage[shortlabels]{enumitem}
\usepackage{hyperref}
\usepackage{version}
\usepackage{bbold}
\usepackage{fontawesome5}
\usepackage{fancyhdr}
\lfoot{}
\usepackage{natbib}
\usepackage{graphicx}
\usepackage{pdfpages}
\usepackage{authblk}
    
\usepackage{blindtext}
\usepackage{multicol}
\usepackage{collcell}
\usepackage{wrapfig}
\usepackage{blkarray}
\usepackage{cancel}
\usepackage{appendix}
\usepackage{float}
\makeatletter
\input{tex/assets/slashbox.sty}
\makeatother
\hypersetup{unicode,colorlinks,citecolor=red,filecolor=black,
  linkcolor=red,urlcolor=cyan}
\setlength{\columnsep}{1cm}
\pagenumbering{gobble}

\def\numtext#1{%
  \ifcase#1\or one\or two\or three\or four\or five\or six\or seven\or eight\or nine\or ten\or
  eleven\or twelve\or thirteen\or fourteen\or fifteen\or sixteen\or seventeen\or eighteen\or nineteen\or twenty\or Lots\fi}

\newcommand{\Byte}[1]{%
  \def\nodecontents{}%
  \foreach[count=\l] \k in {#1} {\xappto{\nodecontents}{\noexpand\nodepart{\numtext{\l}}\k}}%
  \node[rectangle split,rectangle split parts=8,rectangle split horizontal,draw] at (2,2){\nodecontents};}

\newcommand{\D}{{\mathrm{d}}}
\newcommand{\cE}{{\mathcal{E}}}
\newcommand{\cG}{{\mathcal{G}}}
\newcommand{\cX}{{\mathcal{X}}}
\newcommand{\cY}{{\mathcal{Y}}}
\newcommand{\cB}{{\mathcal{B}}}
\newcommand{\cF}{{\mathcal{F}}}
\newcommand{\cO}{{\mathcal{O}}}
\newcommand{\cS}{{\mathcal{S}}}
\newcommand{\cT}{{\mathcal{T}}}
\newcommand{\given}{\, \big| \,}
\newcommand{\cL}{{\mathscr{L}}}
\newcommand{\cD}{{\mathcal{D}}}
\newcommand{\cA}{{\mathscr{A}}}
\newcommand{\cP}{{\mathscr{P}}}
\newcommand{\cR}{{\mathscr{R}}}
\newcommand{\A}{{\mathcal A}}
\newcommand{\one}{\mathbb{1}}
\newcommand{\mathbbm}[1]{\mathbb{#1}}

\providecommand{\C}{}
\renewcommand{\C}{\mathbb C}
\newcommand{\R}{\mathbb R}
\newcommand{\Q}{\mathbb Q}
\newcommand{\N}{\mathbb N}
\newcommand{\T}{\mathbb T}
\newcommand{\F}{\mathbb F}
\renewcommand{\P}{\mathbb P}
\newcommand{\Z}{\mathbb Z}
\newcommand{\E}{\mathbb {E}}
\newcommand{\Pbb}[1]{\P\left( #1\right)}
\newcommand{\Ebb}[1]{\E\left[ #1\right]}

\newcommand{\lb}{\left (}
\newcommand{\rb}{\right )}
\newcommand{\lbb}{\left [}
\newcommand{\rbb}{\right ]}
\newcommand{\labs}{\left |}
\newcommand{\rabs}{\right |}
\newcommand{\lbrb}[1]{\lb #1 \rb}
\newcommand{\lbbrbb}[1]{\lbb#1\rbb}
\newcommand{\labsrabs}[1]{\labs#1\rabs}
\newcommand{\lbbrb}[1]{\lbb#1\rb}
\newcommand{\lbrbb}[1]{\lb#1\rbb}
\newcommand{\langlerangle}[1]{\langle#1\rangle}
\newcommand{\lbcurly}{\{}
\newcommand{\rbcurly}{\}}
\newcommand{\curly}[1]{\lbcurly#1\rbcurly}
\newcommand{\abs}{\labsrabs}
\newcommand{\ind}[1]{1_{#1}}
\newcommand{\indep}{\perp \!\!\! \perp}
\renewcommand{\d}{\, \mathrm{d}}

\newtheorem{theorem}{Theorem}
\newtheorem{problemCnt}{Lemma}
\theoremstyle{definition}
\newtheorem{problem}[problemCnt]{Problem}
\newtheorem{problembook}[problemCnt]{\faBook \,  Problem}
\newtheorem{problempencil}[problemCnt]{\faPencil* \,  Problem}
\newtheorem{problemstar}[problemCnt]{\faStar\,  Problem}
\newtheorem{problemstarpencil}[problemCnt]{\faStar \, \faPencil* \,   Problem}
\newtheorem{probleminteresting}[problemCnt]{\faSmile[regular] \,   Problem}
\newtheorem{solution}{Solution}

\newcommand{\fwcell}[1]{\makebox[\arraystretch\normalbaselineskip]{$#1$}}
\newenvironment{enumeratei}{\begin{enumerate}[(i)]}{\end{enumerate}}
\newenvironment{enumeratea}{\begin{enumerate}[(a)]}{\end{enumerate}}
\newenvironment{enumeraten}{\begin{enumerate}[\upshape 1.]}{\end{enumerate}}

% Compatibility helpers needed by a few files in the combined archive.
\AtBeginDocument{\noindent{\footnotesize English translation prepared with ChatGPT (OpenAI, GPT-6), 26 October 2026, for translation of the Bulgarian text into English and LaTeX formatting.
The Bulgarian originals remain the reference; translation may contain errors.}\par\smallskip}
\newcommand{\MissingGraphic}[1]{%
  \begin{center}\fbox{\parbox{0.82\linewidth}{\centering Missing source image: \texttt{#1}}}\end{center}}



\setlength{\emergencystretch}{3em}
\allowdisplaybreaks
\begin{document}


\noindent Probability and Statistics, Software Engineering, Year 3, 2020/2021 \hfill Written exam, 27.01.2021

\hrulefill

\begin{it}
Your grade is 2 plus the number of points earned. Time allowed: 3 hours. Good luck.

Throughout, we work on a probability space $(\Omega, \cF, \P)$.
\end{it}

\begin{problem} (1 point)
An urn contains balls numbered from 1 to 10 inclusive.
Three balls are drawn without replacement. Define
the following events concerning the balls drawn:
\begin{enumerate}
    \item $A = $\textit{the largest number is 6};
    \item $B = $\textit{at least two balls have even numbers};
    \item $C = $\textit{the sum of the numbers is divisible by 3}.
    \end{enumerate}
\end{problem}
Find the probabilities of these events and
determine whether any are independent. What is $\P(A|B)$?
What is the expected number of even-numbered balls drawn
if sampling is with replacement?

$\bullet$

Let $c=\binom{10}{3}=120$. Then

\[ \P(A) = \frac{\binom{5}{2}}{c} = \frac{1}{12}, \quad \P(B) = \frac{\binom{5}{3} + 5\binom{5}{2}}{c} = \frac12 \, .
\]
For $\P(C)$, either enumerate the combinations directly,
or observe that the possible combinations modulo $3$ are
three equal residues or $0+1+2$. Hence
\[
\P(C) = \frac{3\cdot 4 \cdot 3 + 2 \cdot \binom{3}{3} + \binom{4}{3}}{c}
= \frac{7}{20}\, .
\]
Direct calculation/enumeration gives
\[
\P(A \cap B) = \frac{\binom{2}{2} + \binom{3}{1}\cdot 2}{c} = \frac{7}{120} , \P(A \cap C) = \frac{3}{120} , \P(B \cap C) = \frac{21}{120}\, .
\]
No pair is independent, and
\[\P(A|B) = \frac{\P(A \cap B)}{\P(B)} = \frac{7}{60} \, .\]
With replacement, the number of even-numbered balls drawn
has distribution $Bin(3,1/2)$, so the required expectation
is $3/2$.

\begin{problem}(1 point)
Two instruments measure the concentration of dust particles
in the air.
Suppose the true concentration is $x$ $g/m^3$.
The reading of the first instrument
is normally distributed with mean $x$ and standard deviation ($\sigma$) $0.05x$.
The second reading is also normally distributed,
with mean $x$ and standard deviation
$0.1x$. Which instrument would you use?
For each instrument, what is the probability of an error
greater than $0.1x$?

Someone decides to use the arithmetic mean of the two readings.
If the measurements are independent, what is the probability
of an error greater than $0.1x$ with this procedure?
\end{problem}


$\bullet$
The instrument with smaller variance is preferable. Denote the readings
of the first and second instruments by $\xi_1\sim N(x,(0.05x)^2)$ and
$\xi_2\sim N(x,(0.1x)^2)$. The required probabilities are
\begin{align*}
\P(|\xi_1 - x| > 0.1x) &= \P(\xi_1 < 0.9x) + \P(\xi_1 > 1.1x)
= 2\P(\xi_1 < 0.9x) \\&= 2 \P\left(N(0,1) < \frac{0.9x - x}{0.1x}\right)
 = 2\P(N(0,1) < -1) \\
 &\approx
 0.3174 \, .
\end{align*}
and similarly
\[
\P(|\xi_2 - x| > 0.1x) = 2 \P(N(0,1) < -2) \approx 0.0456
\, .
\]

If we define $\xi_3=1/2(\xi_1+\xi_2)$, then
$\xi_3\sim N(x,1/4(0.05^2+0.1^2)x^2)$ and
\[
\P(|\xi_3 - x| > 0.1x) = 2  \P\left(N(0,1) < \frac{-0.1x}{
\sqrt{1/4( 0.05^2 + 0.1^2)x^2 }}\right)\approx
2\P(N(0,1) < -1.79) \approx 0.0734 \, .
\]



\begin{problem} (1 point)
Let $\xi$ and $\eta$ be independent random variables,
$\xi\sim Exp(2)$ and $\eta\sim U(0,3)$, i.e.
\[
f_\xi(x) = \begin{cases}
2 e^{-2x} & \text{if } x >0\\
0 & \text{otherwise}
\end{cases}; \quad
f_\eta(x) = \begin{cases}
\frac13 & \text{if } 0<x <3\\
0 & \text{otherwise.}
\end{cases}\]
Find the correlation of $\xi$ and $\eta$, $P(\xi<\eta)$ and the density
of $\xi/\eta$.

\end{problem}

$\bullet$

The correlation is 0 because $\xi$ and $\eta$ are independent. Considering
$0<\eta<3$ and $0<\xi<\eta$, we obtain
\[
\P(\xi < \eta) =
\int_0^3\int_0^y 2e^{-2x}\cdot\frac13 \d x \d y
=\frac16\left(5 + \frac{1}{e^6}\right) \approx 0.8334 \, .
\]

Note that independence gives
$f_{\xi, \eta}(x,y) = 2/3 \cdot e^{-2x} \mathbb{1}_{x > 0, 0 < y < 3}$.
Define $X=\xi/\eta$ and $Y=\eta$. The inverse transformation
is
$\xi=XY$ and $\eta=Y$, with Jacobian $Y$. Therefore
\[
f_{X,Y}(x,y) = yf_{\xi,\eta}(xy,y) = y\frac23 e^{-2xy}\mathbb{1}_{xy>0,
0<y<3} \, .
\]
Thus, for $x>0$,
\begin{align*}
f_X(x) &= \int_\infty^\infty f_{X,Y}(x,y) \d y =
\frac23 \int_0^3 ye^{-2xy} \d y
= \frac{1}{6x^2} \int_0^3 -2xy \cdot e^{-2xy} \d (-2xy)\\
&= \frac{1}{6x^2} \int_0^{-6x} t \cdot e^t \d t
 = \frac{e^{-6x}(-6x - 1) + 1}{6x^2} \, ,
\end{align*}
as required.
The last result also gives another way to calculate
the probability requested earlier:
\[
\P(\xi < \eta) = \P\left(\frac{\xi}{\eta}< 1\right) = \int_0^1 f_X(x) \d x
= \int_0^1\frac{e^{-6x}(-6x - 1) + 1}{6x^2} \d x
\]
\begin{problem} (1 point)
Let $\xi$ and $\eta$ have joint density
\[
f_{\xi, \eta}(x, y) = \begin{cases}
c x^3 y^2 & \text{if } 0<x,y<2\\
0 & \text{otherwise.}
\end{cases}
\]
Find the constant $c$, the correlation of $\xi$ and $\eta$, and
$\E(\xi| \eta = 1)$.
\end{problem}
$\bullet$
\[
1 = c \int_0^2 \int_0^2 x^3 y^2 \d x \d y = c \frac{32}{3} \, .
\]
Hence $c=3/32=0.09735$. For $x\in(0,2)$,

\begin{align*}
f_\xi(x) = \int_0^2 cx^3 y^2 \d y = cx^3 \frac{8}{3} = \frac14 x^3 \, .
\end{align*}
Therefore
\[ \E \xi = \frac14 \int_0^2 x^4 \d x = \frac85 = 1.6, \quad
D \xi = \frac14\int_0^2 x^5 \d x - \frac{64}{25} = \frac{8}{75}\approx 0.1067 \, .\]
Similarly, for $y\in(0,2)$,
\[
f_\eta(y) = c\int_0^2 x^3 y^2\d x = \frac{3}{8}y^2\, ,
\]
\[
\E \eta = \frac38\int_0^2 y^3 \d y  =
\frac{3}{2} =1.5 , \quad
D \eta = \frac38\int_0^2 y^4 \d y - \frac94 = \frac{3}{20} = 0.15\, \, .
\]
Next,
\[
\E \xi \eta = \frac{3}{32} \int_0^2 \int_0^2 xy(x^3y^2) \d x \d y
= \frac{12}{5} = 2.4
\]
Therefore
\[
Cor(\xi, \eta) =\frac{ \E \xi\eta - \E \xi \E \eta}
{\sqrt{D\xi D\eta}}= 0 \, .
\]
For the last question,
\begin{align*}
\E (\xi | \eta = 0.5) &= \int_0^2 x f_{\xi|\eta}(x | y = 1) \d x
= \int_0^2 x \frac{f_{\xi, \eta}(x, 1)}{f_\eta(1)} \d x
= \int_0 ^2 x \frac{x^3}{4}\d x
= \frac85=1.6 \, .
\end{align*}

\end{document}
